% =====================================================================
% Frente 3 --- Centralidades e desempenho (Iago)
% =====================================================================
\section{Frente 3 --- Centralidades e desempenho}

\noindent\textbf{Responsável:} Iago.
\textbf{No relatório:} centralidades; avaliação de desempenho.
\textbf{Na Etapa 1:} diagrama de sequência.

\subsection{O que foi implementado}

{\small\begin{tabularx}{\linewidth}{@{}L{3.9cm} L{4.7cm} Y L{1.9cm}@{}}
\toprule
Função & Arquivo & Complexidade & Situação \\
\midrule
\codigo{intermediacao} (exata) & \codigo{metricas/centralidade.py} & $O(nm)$~\cite{brandes2001} & \feito \\
\codigo{intermediacao} ($k$ fontes) & \codigo{metricas/centralidade.py} & $O(km)$~\cite{brandes2007} & \feito \\
\codigo{pagerank} & \codigo{metricas/centralidade.py} & $O(m)$ por iteração~\cite{brin1998} & \feito \\
\codigo{numero\_de\_nucleo}, \codigo{k\_nucleo} & \codigo{metricas/nucleo\_k.py} & $O(m)$~\cite{batagelj2003} & \feito \\
\codigo{medir.py} (tempo, memória, TEPS) & \codigo{desempenho/} & --- & \afazer \\
\codigo{comparar\_networkx.py} & \codigo{desempenho/} & --- & \afazer \\
\bottomrule
\end{tabularx}}

\subsection{Decisões e justificativas}
\begin{itemize}[leftmargin=*]
  \item Listas indexadas $0..n-1$ (\codigo{IndexadorDeNos}) no Brandes, no PageRank e no k-core: mais rápidas que os dicionários do \codigo{Grafo}.
  \item No PageRank, nós sem saída repartem o valor entre todos os nós, e a aresta é escolhida em proporção ao peso.
  \item Se o PageRank não converge em \codigo{max\_iteracoes} (acontece em grafos bipartidos, como a estrela, com a tolerância $10^{-8}$), a função emite um aviso e devolve a última iteração, com erro da ordem de $0{,}85^{100} \approx 10^{-7}$, em vez de lançar exceção.
  \item No k-core, laços são ignorados; no direcionado, o grau é entrada + saída, como no NetworkX.
\end{itemize}

\subsection{Corretude: testes e resultados}

{\small\begin{tabularx}{\linewidth}{@{}Y L{2.5cm} L{2.2cm} L{3.8cm}@{}}
\toprule
Teste & Base & Obtido & Referência \\
\midrule
Betweenness normalizada dos nós 0, 33 e 32 & Karate & 0,4376; 0,3041; 0,1452 & 0,4376; 0,3041; 0,1452 \\
PageRank ($d = 0{,}85$) dos nós 33, 0 e 32; soma & Karate & 0,1009; 0,097; 0,0717; 1 & 0,1009; 0,097; 0,0717; 1 \\
Núcleo máximo & Karate & 4 (os mesmos 10 nós) & 4 (nós 0, 1, 2, 3, 7, 8, 13, 30, 32, 33) \\
Betweenness e PageRank máximos & Twitch PT & \falta & nó 127 (0,0993 e 0,0119) \\
Núcleo máximo & Twitch PT & \falta & 37, com 154 nós \\
Centro da estrela; grafo completo & sintético & $(n-1)(n-2)/2$; 0 & $(n-1)(n-2)/2$; 0 \\
\bottomrule
\end{tabularx}}

\subsection{Desempenho}
\begin{itemize}[leftmargin=*]
  \item Ambiente: \pendente{CPU, RAM, sistema, versão do Python e do NetworkX}
  \item Metodologia: mediana de 5 execuções com \codigo{time.perf\_counter}; pico de memória com
        \codigo{tracemalloc}; grafos ER, BA e WS de $10^3$ a $10^5$ nós e as 4 bases.
  \item Resultados: \pendente{tabela nosso código vs. NetworkX; inclinação do log-log}
\end{itemize}

\subsection{Figuras e tabelas para o relatório}
\begin{itemize}[leftmargin=*]
  \item \pendente{gráfico log-log de tempo $\times$ $(n+m)$, com retas $O(n+m)$ e $O(nm)$}
\end{itemize}

\subsection{Limitações e pendências}
\begin{itemize}[leftmargin=*]
  \item \pendente{ex.: betweenness exata inviável em grafos grandes; usar amostragem~\cite{brandes2007}}
\end{itemize}

\subsection{Registro}
\begin{description}[leftmargin=2.6cm,style=nextline]
  \registro{2026-11-13}{Diagrama de sequência da Etapa 1 e intermediação de Brandes (exata e por $k$ origens), PageRank e k-core, com testes.}
\end{description}
